% Proposal paragraphs retained; scope and comparison clarification added.
\subsection{Research Design and Methodological Rationale}

We propose a controlled, quantitative 2D simulation of two unmanned surface vehicles (USVs) surveying a known area under intermittent direct radio communication. We compare reactive online resynchronization with prediction-aware pre-outage task reservation; a fixed offline partition is a communication-independent reference. Channel-fade persistence (target mean fade duration) is the principal varying condition. Mission completion time is the primary outcome, with redundant coverage and communication cost as secondary measures.

The three online policies---reactive, non-predictive long backlog, and prediction-triggered reservation---will share the motion and sensing models, task allocator, pre-generated fade and slowdown traces, update interval, and maximum message budget. They will differ in how they select the reservation capacity, including whether an outage warning is used. Prediction-triggered reservation will be evaluated with both persistence and trend predictors, giving five experimental configurations in total when the fixed partition is included. Paired simulations can therefore test whether anticipating an outage changes allocations and outcomes. Real-boat experiments could strengthen external validity, but are infeasible within the available resources and make link conditions harder to control. An optional distance-loss sensitivity check will probe the simplified model after the core comparisons are complete. Software tests verify implementation, not maritime realism; conclusions will be limited to the modeled two-USV scenario rather than real-world navigation safety, radio performance, or seabed-map accuracy.

The implementation will start with deterministic connected coverage, add a single slowdown and controlled outage, and then proceed to randomized paired comparisons. A non-predictive long-backlog control is added to distinguish a useful forecast from simply having more work reserved. Finite-speed motion and an explicit turning cost are needed for task timing; hydrodynamics, currents, collision avoidance, bathymetry, acoustic links, and measured sea-channel validation remain outside the core study. These omissions limit external validity, rather than being evidence that realism does not matter.

We will implement established prediction components ourselves: persistence (the naive last-observation forecast) \cite{hyndman2021forecasting} and ordinary least-squares linear extrapolation \cite{nistLeastSquares}. We are not reusing a published maritime prediction algorithm or claiming a new forecasting method. The signal-motion predictor combines a constant-relative-velocity distance forecast with a fitted signal residual; that combination and the outage-warning trigger are specified for this study in Section 2.3. The contribution being evaluated is the coordination benefit and resource trade-off of using these simple warnings under the same reservation action, not an improvement in state-of-the-art marine channel prediction.
